%% Table 1: Capability--Threat Mapping (from §2)
%% Requires: \usepackage{tabularx}

\begin{table}[htbp]
\centering
\caption{Mapping generative AI capabilities to security threats.}
\label{tab:cap-threat}
\tiny
\setlength{\tabcolsep}{3pt}
\renewcommand{\arraystretch}{0.88}
\begin{tabularx}{\columnwidth}{l l X c}
\toprule
\textbf{Capability} & \textbf{Enabling Property} & \textbf{Threat Vectors Enabled} & \textbf{\S} \\
\midrule
High-fidelity generation & Cross-modal human parity & Undetectable synthetic media; deepfake impersonation; evidence fabrication & \ref{sec:threat-deepfake} \\
Low-cost production & Open weights; sub-\$0.30/M-token APIs & Automated phishing at scale; commodity malware generation & \ref{sec:threat-phishing} \\
Controllability & Fine-tuning (LoRA); few-shot adaptation & Targeted identity cloning; stylometric forgery; personalized social engineering & \ref{sec:threat-ncii} \\
Reasoning \& long context & Chain-of-thought; 1M+ token windows & Multi-step attack planning; codebase-aware exploitation & \ref{sec:threat-model} \\
Autonomous agency & Tool invocation; multi-step execution & Autonomous system compromise; persistent unauthorized actions & \ref{sec:threat-ipi} \\
Tool ecosystem (MCP) & Standardized agent--tool interface & Tool poisoning; cross-server composition attacks; supply chain exploitation & \ref{sec:threat-adversarial} \\
\bottomrule
\end{tabularx}
\end{table}
